% Generated draft: requires mathematical and source review.
\documentclass[11pt]{article}
\usepackage[margin=25mm]{geometry}
\usepackage{fontspec}
\setmainfont{Latin Modern Roman}
\usepackage{amsmath,amssymb,mathtools,unicode-math}
\setmathfont{Latin Modern Math}
\usepackage{enumitem,needspace,xurl,hyperref}
\setlength{\parindent}{0pt}
\setlength{\parskip}{5pt}
\setlength{\emergencystretch}{3em}
\newcommand{\N}{\mathbb{N}}
\newcommand{\Z}{\mathbb{Z}}
\newcommand{\Q}{\mathbb{Q}}
\newcommand{\R}{\mathbb{R}}
\newcommand{\C}{\mathbb{C}}
\newcommand{\F}{\mathbb{F}}
\newcommand{\A}{\mathbb{A}}
\newcommand{\PP}{\mathbb{P}}
\newcommand{\E}{\mathbb{E}}
\newcommand{\eps}{\varepsilon}
\newcommand{\dd}{\,\mathrm{d}}
\begin{document}
\section*{UnsolvedMath ID: 20003147}
\textbf{Source problem code:} catalogue entry
\subsection{Definitions and mathematical statement}
The source asks whether a manifold admitting \(\sec>0\), \(\sec\geq0\), or almost nonnegative sectional curvature must admit a nontrivial smooth circle action. For a closed connected smooth manifold \(M\), define its symmetry degree by
\[
 \operatorname{symdeg}(M)=\sup\{\dim G:G\text{ is a compact Lie group acting smoothly and effectively on }M\}.
\]
Positive symmetry degree is equivalent to the existence of an effective circle action; a nontrivial circle action can be made effective by dividing out its finite kernel. The question concerns \(\operatorname{Diff}(M)\), not just the isometry group of one chosen curvature metric. The source does not explicitly impose compactness or simple connectivity.

The supplied exact flat-manifold criterion is as follows. For every closed connected flat manifold
\[
 F=\R^n/\Gamma,
\]
with Bieberbach group \(\Gamma\), translation lattice, and finite linear holonomy \(H\), the aid states
\[
 \operatorname{symdeg}(F)
 =\operatorname{rank}Z(\Gamma)
 =\dim(\R^n)^H
 =b_1(F).
\]
Here \(Z(\Gamma)\) is the center, rank is its free abelian rank, \((\R^n)^H\) is the holonomy-fixed vector space, and \(b_1(F)=\dim_\Q H_1(F;\Q)\). Consequently,
\[
 F\text{ admits a nontrivial smooth circle action}
 \Longleftrightarrow b_1(F)>0.
\]
The formula is reported for arbitrary effective smooth compact-group actions on the underlying flat manifold, not only affine actions or isometries of its flat metric.

The aid's counterexample is the Hantzsche--Wendt closed flat three-manifold, with
\[
 H_1(F;\Z)\cong(\Z/4\Z)^2,
 \qquad b_1(F)=0,
 \qquad\operatorname{symdeg}(F)=0.
\]
It has \(\sec=0\), so it refutes the literal nonnegative-curvature branch. It also satisfies the standard almost-nonnegative convention of metrics \(g_j\) with \(\sec(g_j)\geq-1/j\), \(\operatorname{diam}(g_j)\leq1\), after a fixed rescaling of the flat metric. Thus it also refutes that branch under this convention.

This example is not simply connected and does not have strictly positive curvature. The aid expressly says the official AIM wording omitted the simply-connected hypothesis in Grove's underlying formulation. It does not use this flat example to refute the strictly positive branch or the simply connected variant. The registry's solved/refuted labels concern the literal unrestricted assertion in the stated nonnegative settings, rather than a classification of all three strengthened versions.
\subsection{Short English statement}
The supplied flat criterion makes symmetry degree equal to the first Betti number. The Hantzsche--Wendt manifold therefore has no smooth circle action despite zero curvature. This refutes the literal nonnegative and almost-nonnegative claims, while leaving the strictly positive and simply connected versions outside this counterexample.
\subsection{Sources}
\begin{itemize}
\item Local UnsolvedMath registry, permanent ID \(20003147\), exact source code \texttt{catalogue entry}; canonical \texttt{aim-topology-notes.json}, \texttt{notes[234]}, source tag \texttt{problem}.
\item AIM workshop list ``Open Problems in Non-Negative Sectional Curvature,'' workshop ``Manifolds with nonnegative sectional curvature,'' Problem 22: \url{https://aimath.org/WWN/nnsectcurvature/nnsectcurvature.pdf}. The supplied verification calls this wording exact; the PDF was not fetched.
\item The aid names the Hantzsche--Wendt manifold and refers to Grove's underlying simply connected formulation, but gives no full authoritative references for that formulation, the smooth-action symmetry theorem, or the homology computation. The original literature field is empty.
\item Research attempt 1: result \texttt{counterexample}, problem status at run \texttt{refuted}, rigor label \texttt{proved}, novelty confidence \texttt{low}; classification \texttt{SOLVED-IN-LITERATURE}, registry status \texttt{solved}. Created and updated 14 August 2026. These are supplied assessments for the literal formulation, not independent acceptance of stronger claims.
\item Attribution: Ulam AI / UnsolvedMath Contributors, CC BY 4.0. The machine-generated research aid is not peer reviewed and requires independent verification before citation.
\end{itemize}
\end{document}
